\chapter{Overview of the Dichotomies}
\label{ch:dichotomy-overview}

We identify twenty-four fundamental \emph{computational dichotomies}
--- pairs of implementation strategies that compute the same function.
Each dichotomy is a path constructor in the HIT $\mathsf{PyComp}$
(Definition~\ref{def:pycomp}).

The dichotomies are organized into four groups:

\begin{center}
\begin{tabular}{c l l}
  \textbf{Group} & \textbf{Dichotomies} & \textbf{OTerm Types Involved} \\
  \hline
  \emph{Control Flow} & D1--D8 & \texttt{OFix}, \texttt{OFold}, \texttt{OCase}, \texttt{OLam} \\
  \emph{Data Structure} & D9--D14 & \texttt{OOp}, \texttt{OFold}, \texttt{ODict} \\
  \emph{Algorithm Strategy} & D15--D20 & \texttt{OFix}, \texttt{OFold}, \texttt{OAbstract} \\
  \emph{Language Feature} & D21--D24 & \texttt{OCatch}, \texttt{OCase}, \texttt{OLam} \\
\end{tabular}
\end{center}

Each dichotomy manifests in two complementary ways within the system:
\begin{enumerate}
  \item \textbf{Normalizer rewrite} (in \texttt{denotation.py}, phases 1--7):
    a directed rule that reduces one form to a canonical OTerm,
    applied during the 7-phase normalization pass.
  \item \textbf{Path constructor} (in \texttt{path\_search.py}):
    a bidirectional axiom that, given an OTerm matching one paradigm,
    produces the equivalent OTerm in the other paradigm together with
    a named proof step.
\end{enumerate}
When both sides of a dichotomy normalize to the same canonical OTerm
the path is $\mathsf{refl}$.  When normalization produces distinct
but equivalent OTerms, the HIT structural prover
(\texttt{hit\_path\_equiv}) applies the path constructor to close the
gap.

\section{Complete Dichotomy Table}

\begin{longtable}{r p{3.5cm} p{3.5cm} p{3.0cm} p{2.2cm}}
  \toprule
  \textbf{\#} & \textbf{Dichotomy} & \textbf{Path Constructor} &
  \textbf{Normalizer Phase} & \textbf{Axiom fn} \\
  \midrule
  \endhead
  \multicolumn{5}{l}{\emph{Control Flow Dichotomies}} \\
  \midrule
  D1 & Recursive $\leftrightarrow$ Iterative &
    Nat-eliminator uniqueness &
    Phase 3 (fold) &
    \texttt{\_axiom\_d1} \\
  D2 & Nested helpers $\leftrightarrow$ Flat &
    $\beta$-reduction &
    Phase 1 ($\beta$) &
    \texttt{\_axiom\_d2} \\
  D3 & \texttt{while} $\leftrightarrow$ \texttt{for} &
    Loop normal form &
    Phase 3 (fold) &
    --- \\
  D4 & Comprehension $\leftrightarrow$ Loop &
    Catamorphism fusion &
    Phase 4 (HOF) &
    \texttt{\_axiom\_d4} \\
  D5 & \texttt{reduce} $\leftrightarrow$ Loop &
    Fold universality &
    Phase 3 (fold) &
    \texttt{\_axiom\_d5} \\
  D6 & Generator $\leftrightarrow$ Eager &
    Laziness adjunction &
    Phase 4 (HOF) &
    --- \\
  D7 & Tail-recursive $\leftrightarrow$ Loop &
    CPS transform &
    Phase 3 (fold) &
    --- \\
  D8 & Multi-return $\leftrightarrow$ Single &
    Section merging &
    Phase 7 (piecewise) &
    \texttt{\_axiom\_d8} \\
  \midrule
  \multicolumn{5}{l}{\emph{Data Structure Dichotomies}} \\
  \midrule
  D9 & Stack $\leftrightarrow$ Recursion &
    Defunctionalization &
    Phase 5 (unify) &
    --- \\
  D10 & \texttt{heapq} $\leftrightarrow$ \texttt{bisect} &
    Priority queue abstr. &
    Phase 5 (unify) &
    --- \\
  D11 & \texttt{OrderedDict} $\leftrightarrow$ LinkedList &
    ADT refinement &
    Phase 5 (unify) &
    --- \\
  D12 & \texttt{defaultdict} $\leftrightarrow$ \texttt{setdefault} &
    Map construction equiv &
    Phase 5 (unify) &
    --- \\
  D13 & \texttt{Counter} $\leftrightarrow$ Manual count &
    Histogram equivalence &
    Phase 5 (unify) &
    \texttt{\_axiom\_d13} \\
  D14 & Array $\leftrightarrow$ Dict table &
    Indexing abstraction &
    Phase 5 (unify) &
    --- \\
  \midrule
  \multicolumn{5}{l}{\emph{Algorithm Strategy Dichotomies}} \\
  \midrule
  D15 & BFS $\leftrightarrow$ DFS &
    Traversal-order invariance &
    Phase 6 (quotient) &
    --- \\
  D16 & Memoized $\leftrightarrow$ DP &
    Eval-order independence &
    Phase 5 (unify) &
    \texttt{\_axiom\_d16} \\
  D17 & Divide-and-conquer $\leftrightarrow$ Iterative &
    Tree fold = list fold &
    Phase 4 (HOF) &
    \texttt{\_axiom\_d17} \\
  D18 & Greedy $\leftrightarrow$ DP &
    Matroid/greedy theorem &
    --- &
    --- \\
  D19 & Sort-then-scan $\leftrightarrow$ Sweep &
    Sorting abstraction &
    Phase 5 (unify) &
    \texttt{\_axiom\_d19} \\
  D20 & Different algorithms for same spec &
    Specification HIT &
    --- &
    \texttt{\_axiom\_d20} \\
  \midrule
  \multicolumn{5}{l}{\emph{Language Feature Dichotomies}} \\
  \midrule
  D21 & \texttt{isinstance} $\leftrightarrow$ Dispatch &
    Fiber projection &
    Phase 7 (piecewise) &
    --- \\
  D22 & \texttt{try/except} $\leftrightarrow$ Conditional &
    Effect elimination &
    Phase 2 (ring) &
    \texttt{\_axiom\_d22} \\
  D23 & Sorted-then-process $\leftrightarrow$ Direct &
    Canonical ordering path &
    Phase 5 (unify) &
    --- \\
  D24 & Lambda $\leftrightarrow$ Named function &
    $\eta$-expansion &
    Phase 1 ($\beta$) &
    \texttt{\_axiom\_d24} \\
  \bottomrule
\end{longtable}

\begin{remark}[Axiom registration]
All root axiom functions listed above are collected in the list
\texttt{\_ROOT\_AXIOMS} (path\_search.py, line 1580).
During bidirectional BFS, each axiom is applied to every sub-expression
of the current OTerm; the resulting rewrites form the neighbour set
for the next BFS layer.  Axioms marked ``---'' are handled entirely
by the normalizer and never need a separate path-search step.
\end{remark}


\chapter{Control Flow Dichotomies (D1--D8)}
\label{ch:cf-dichotomies}

Each section below follows the same structure: (1)~two concrete programs
illustrating paradigm A and paradigm B, (2)~how each is compiled to
OTerms, (3)~the formal path constructor definition, (4)~a soundness
theorem with proof, (5)~the implementation mapping to
\texttt{path\_search.py} and \texttt{denotation.py}, and (6)~a benchmark
example where this dichotomy applies.

%% ────────────────────────────────────────────────────
\section{D1: Recursive $\leftrightarrow$ Iterative}
\label{sec:d1}
%% ────────────────────────────────────────────────────

\textbf{Paradigm A} --- recursive descent:
\begin{lstlisting}
def f(n):
    if n <= 1: return 1
    return n * f(n - 1)
\end{lstlisting}

\textbf{Paradigm B} --- iterative accumulation:
\begin{lstlisting}
def g(n):
    result = 1
    for i in range(2, n + 1):
        result *= i
    return result
\end{lstlisting}

\textbf{Compilation of A.}
The compiler sees a self-recursive function with one parameter.  It
invokes the \emph{Nat-eliminator extractor}, which detects:
\begin{itemize}
  \item Base guard: \texttt{n <= 1}, so $\mathsf{threshold} = 1$.
  \item Base value: literal \texttt{1}, so $\mathsf{base} = 1$.
  \item Step: \texttt{n * f(n-1)}, a $\mathsf{BinOp}(\times)$ with
    the parameter on the left and the recursive call on the right.
    The operation tag is $\mathsf{MUL}$.
\end{itemize}
The extractor produces: $\den{f} = \fold(\mathsf{MUL}, 1,
\mathsf{ival}(p_0) + 1, \mathsf{IntObj}(1))$.

\textbf{Compilation of B.}
The compiler sees an assignment \texttt{result = 1} (init), then
\texttt{for i in range(2, n+1): result *= i}.  The loop-to-fold
optimizer (Phase~3, \texttt{\_phase3\_fold}) detects:
\begin{itemize}
  \item Single accumulator \texttt{result} with init $= 1$.
  \item \texttt{AugAssign} with $\mathsf{Mult}$: \texttt{result *= i}.
    The function \texttt{\_detect\_augassign\_fold} returns semantic
    key \texttt{imul}.
  \item $\mathsf{range}(2, n+1)$ gives $\mathsf{start} = 2$,
    $\mathsf{stop} = \mathsf{ival}(p_0) + 1$.
\end{itemize}
The optimizer produces: $\den{g} = \fold(\mathsf{MUL}, 1,
\mathsf{ival}(p_0) + 1, \mathsf{IntObj}(1))$.

\textbf{Path.}
$\den{f}$ and $\den{g}$ are \emph{the same OTerm}.  The path is
$\mathsf{refl}$ (reflexivity) --- no axiom application needed.

When the step has $f(n-1) * n$ instead of $n * f(n-1)$, the terms
differ: $\fold(\mathsf{MUL}, \ldots)$ vs.\ $\fold(\mathsf{MUL},
\ldots)$ but with swapped operands inside the RecFunction body.
The $\mathsf{mul\_comm}$ path constructor bridges the gap:
the ring-normalization phase (Phase~2, \texttt{\_phase2\_ring})
sorts commutative operands into canonical order, yielding
$\mathsf{refl}$ after normalization.

\begin{definition}[D1 Path Constructor --- formal]
\label{def:d1-path}
Define the path constructor for $\mathsf{D1}$ as the family
\[
  \mathsf{D1} \;\coloneqq\; \lambda\, b\, s\, n.\;
  \mathsf{path}(\mathsf{OFix}(\mathsf{h},\;
    \mathsf{case}(\mathsf{le}(n, 0),\; b,\; s(n, \mathsf{h}(n-1)))),\;
  \mathsf{OFold}(s, b, n))
\]
where $b : \mathsf{OTerm}$ is the base value, $s$ is the step
operation name, and $n$ is the collection/range bound.  The path
witnesses that a recursive fixed point with accumulation pattern
is equal to the corresponding fold.
\end{definition}

\begin{theorem}[D1 Soundness]
\label{thm:d1-sound}
For any base value $b$, binary operation $\otimes$ satisfying
$\otimes(x, y) = \otimes(y, x)$ (commutativity is optional;
required only for operand-swapped variants), and natural number $n$:
\[
  \mathsf{D1} : \prod_{b,\, \otimes,\, n}
  \Path\!\big(\rec_\Nat(b, \lambda k.\, k \otimes \rec_\Nat(b, \ldots)(k-1)),\;
  \fold(\otimes, b, n)\big)
\]
\end{theorem}

\begin{proof}
By the uniqueness of the $\Nat$-eliminator (the universal property of
$\Nat$ as an initial algebra):  any function $h : \Nat \to A$ satisfying
$h(0) = b$ and $h(k+1) = k \otimes h(k)$ is uniquely determined.
The recursive form $\rec_\Nat(b, s)$ satisfies this recurrence by
definition.  The fold $\fold(\otimes, b, n)$ unfolds to the same
recurrence: $\fold(\otimes, b, 0) = b$ and $\fold(\otimes, b, k+1)
= \otimes(k, \fold(\otimes, b, k))$.  By initiality, both are the
same morphism from $\Nat$ to $A$, hence equal.
\end{proof}

\textbf{Implementation mapping.}
\begin{itemize}
  \item \textbf{Normalizer}: \texttt{\_phase3\_fold} in
    \texttt{denotation.py} (line 1379) rewrites builtin reducers to
    \texttt{OFold}; the compiler's \texttt{\_exec\_for\_ot}
    (line 762) detects augmented-assignment loops and emits
    \texttt{OFold} with a semantic key via
    \texttt{\_detect\_augassign\_fold} (line 808).
  \item \textbf{Path search}: \texttt{\_axiom\_d1\_fold\_unfold}
    (line 188) applies in direction OFix $\to$ OFold by calling
    \texttt{\_try\_extract\_fold\_from\_fix}.
  \item \textbf{HIT prover}: \texttt{\_fix\_fold\_equiv}
    (line 1492) handles the cross-type case
    $\mathsf{OFix} \leftrightarrow \mathsf{OFold}$ by comparing
    fix-point names, free variables, and spec identifications.
\end{itemize}

\textbf{Benchmark example.}
The factorial pair (\texttt{demo\_factorial\_rec.py} vs.\
\texttt{demo\_factorial\_iter.py}): the recursive version compiles to
$\mathsf{OFix}[\mathsf{h}](\mathsf{case}(\mathsf{le}(p_0, 1),
\mathsf{1}, \mathsf{mult}(p_0, \mathsf{h}(\mathsf{sub}(p_0,1)))))$
while the iterative version compiles directly to
$\mathsf{OFold}(\mathsf{imul}, \mathsf{1}, p_0)$.
Phase~3 normalizes the recursive form to the same fold,
producing $\mathsf{refl}$.


%% ────────────────────────────────────────────────────
\section{D2: Nested Helpers $\leftrightarrow$ Flat Code}
\label{sec:d2}
%% ────────────────────────────────────────────────────

\textbf{Paradigm A} --- nested helper function:
\begin{lstlisting}
def f(x):
    def helper(a, b):
        return a * a + b * b
    return helper(x, x + 1)
\end{lstlisting}

\textbf{Paradigm B} --- inlined flat code:
\begin{lstlisting}
def g(x):
    return x * x + (x + 1) * (x + 1)
\end{lstlisting}

\textbf{Compilation of A.}
The compiler encounters \texttt{def helper}: it stores an
$\mathsf{OLam}((\texttt{a}, \texttt{b}), \mathsf{body})$ in the
environment (denotation.py, line 241).  When it reaches
\texttt{helper(x, x+1)}, it performs $\beta$-reduction:
substituting $a \mapsto p_0$ and
$b \mapsto \mathsf{add}(p_0, \mathsf{1})$
into the body.  Result:
\[
  \den{f} = \mathsf{add}(\mathsf{mult}(p_0, p_0),\;
    \mathsf{mult}(\mathsf{add}(p_0, \mathsf{1}),
                   \mathsf{add}(p_0, \mathsf{1})))
\]

\textbf{Compilation of B.}
Direct compilation produces the identical OTerm (same expression
tree).

\textbf{Path.}
$\den{f} \equiv \den{g}$ by $\beta$-reduction.  The path is
$\mathsf{refl}$ after inlining.

\begin{definition}[D2 Path Constructor --- formal]
\label{def:d2-path}
\[
  \mathsf{D2} \;\coloneqq\; \lambda\, \vec{x}\, \vec{a}\, e.\;
  \mathsf{path}\!\big(
    \mathsf{OOp}(\mathsf{apply},\,
      \mathsf{OLam}(\vec{x},\, e),\, \vec{a}),\;
    e[\vec{x} \coloneqq \vec{a}]
  \big)
\]
where $e[\vec{x} \coloneqq \vec{a}]$ denotes simultaneous
capture-avoiding substitution of actuals $\vec{a}$ for formals
$\vec{x}$ in the body $e$.
\end{definition}

\begin{theorem}[D2 Soundness --- $\beta$-Reduction]
\label{thm:d2-sound}
For any $\mathsf{OLam}(\vec{x}, e)$ and arguments $\vec{a}$ of
matching arity:
\[
  \mathsf{D2} : \Path\!\big(
    (\lambda \vec{x}.\, e)(\vec{a}),\;
    e[\vec{x} \coloneqq \vec{a}]
  \big)
\]
\end{theorem}

\begin{proof}
This is the $\beta$-rule of the simply typed $\lambda$-calculus,
which is a definitional equality in cubical type theory.  The
substitution $e[\vec{x} \coloneqq \vec{a}]$ is computed by
\texttt{\_subst\_term} (path\_search.py) and is capture-avoiding
because OTerm variables are globally fresh parameter names $p_i$.
\end{proof}

\textbf{Implementation mapping.}
\begin{itemize}
  \item \textbf{Normalizer}: \texttt{\_phase1\_beta} in
    \texttt{denotation.py} (line 1096) collapses trivial cases
    ($\mathsf{case}(\top, t, f) \to t$, $\mathsf{case}(c, x, x) \to x$).
    The compiler itself performs $\beta$-reduction at call sites by
    inlining the stored \texttt{OLam} body.
  \item \textbf{Path search}: \texttt{\_axiom\_d2\_beta}
    (line 221) fires on $\mathsf{OOp}(\mathsf{apply}, \mathsf{OLam}(\ldots), \ldots)$
    nodes.  It pattern-matches the apply-of-lambda form, substitutes
    actuals for formals, and returns the reduced body tagged
    \texttt{D2\_beta}.
\end{itemize}

\textbf{Benchmark example.}
22 benchmark pairs involve nested helper functions.  A typical case:
one program defines \texttt{def helper(a, b): return a + b} then
calls \texttt{helper(x, y)}, while the equivalent program writes
\texttt{x + y} directly.  After $\beta$-reduction both yield
$\mathsf{add}(p_0, p_1)$.


%% ────────────────────────────────────────────────────
\section{D3: \texttt{while} $\leftrightarrow$ \texttt{for}}
\label{sec:d3}
%% ────────────────────────────────────────────────────

\textbf{Paradigm A} --- \texttt{while} with explicit counter:
\begin{lstlisting}
def f(n):
    result, i = 0, 1
    while i <= n:
        result += i
        i += 1
    return result
\end{lstlisting}

\textbf{Paradigm B} --- \texttt{for} over range:
\begin{lstlisting}
def g(n):
    result = 0
    for i in range(1, n + 1):
        result += i
    return result
\end{lstlisting}

\textbf{Compilation of A.}
The compiler encounters a \texttt{while} loop.  It identifies modified
variables $\{\texttt{result}, \texttt{i}\}$ and creates a RecFunction
for each:
\[
  \mathsf{wh\_result}(k) = \begin{cases}
    \mathsf{IntObj}(0) & \text{if } k > 50 \lor \lnot(i \leq n) \\
    \mathsf{add}(\mathsf{wh\_result}(k+1), \ldots)
      & \text{otherwise}
  \end{cases}
\]

\textbf{Compilation of B.}
The compiler encounters \texttt{for i in range(1, n+1)}.  It extracts
$\mathsf{start} = 1$, $\mathsf{stop} = \mathsf{ival}(p_0) + 1$.
The single-accumulator fold pattern
(\texttt{result += i}) is detected by \texttt{\_detect\_augassign\_fold}
which returns semantic key \texttt{iadd}, producing:
\[
  \den{g} = \mathsf{OFold}(\mathsf{iadd},\;
    \mathsf{0},\; p_0)
\]

\textbf{Path.}
Both RecFunctions satisfy the same recurrence:
$h(0) = 0$, $h(k) = h(k-1) + k$.
The \texttt{while} version's guard $i \leq n$ is equivalent to the
\texttt{for} range $[1, n+1)$, and both step by 1.  The RecFunctions
unfold identically.

When Z3 checks $\den{f}(p_0) \neq \den{g}(p_0)$ on the integer fiber
($\mathsf{tag}(p_0) = \mathsf{TInt}$), it finds UNSAT because both
RecFunctions compute the same value for every integer input.

\begin{definition}[D3 Path Constructor --- formal]
\label{def:d3-path}
\[
  \mathsf{D3} \;\coloneqq\; \lambda\, g\, s\, \mathit{init}.\;
  \mathsf{path}\!\big(
    \mathsf{OFix}(\mathsf{h},\;
      \mathsf{case}(g,\; s(\mathsf{h}),\; \mathit{init})),\;
    \mathsf{OFold}(s,\; \mathit{init},\; \mathsf{range}(g))
  \big)
\]
where $\mathsf{range}(g)$ denotes the iteration range induced by
the while-guard $g$ with unit step.
\end{definition}

\textbf{Automatic inference.}
The compiler automatically canonicalizes both loop forms to
RecFunctions.  The \texttt{for}-range version gets the additional
fold optimization.  Z3's recursive function reasoning handles the
comparison.  The path is:
$\mathsf{D3} = \mathsf{while\_rec} \xrightarrow{\text{unfold}} 
\mathsf{for\_fold}$, mediated by the shared recurrence structure.
D3 has no dedicated axiom function in \texttt{path\_search.py}
because the normalizer (phases 1--3) already reduces both loop
forms to the same canonical OTerm.


%% ────────────────────────────────────────────────────
\section{D4: Comprehension $\leftrightarrow$ Loop}
\label{sec:d4}
%% ────────────────────────────────────────────────────

\textbf{Paradigm A} --- list comprehension:
\begin{lstlisting}
def f(xs):
    return [x * 2 for x in xs]
\end{lstlisting}

\textbf{Paradigm B} --- explicit loop with append:
\begin{lstlisting}
def g(xs):
    result = []
    for x in xs:
        result.append(x * 2)
    return result
\end{lstlisting}

\textbf{Compilation of A.}
The compiler encounters a \texttt{ListComp}.  The function
\texttt{\_compile\_comprehension\_ot} (denotation.py, line 556)
builds an $\mathsf{OMap}$:
\[
  \den{f} = \mathsf{OMap}\!\big(
    \mathsf{OLam}(x,\, \mathsf{mult}(x, 2)),\; p_0\big)
\]
The transform is captured as an \texttt{OLam} with the comprehension
variable as the bound parameter and the element expression as the
body.

\textbf{Compilation of B.}
The compiler sees \texttt{result = []} followed by
\texttt{for x in xs: result.append(x * 2)}.  The function
\texttt{\_detect\_map\_pattern} (denotation.py, line 774) recognizes
the append-in-loop idiom and produces the same $\mathsf{OMap}$:
\[
  \den{g} = \mathsf{OMap}\!\big(
    \mathsf{OLam}(x,\, \mathsf{mult}(x, 2)),\; p_0\big)
\]

\textbf{Path.}
Both programs normalize to structurally identical \texttt{OMap} terms.
The path is $\mathsf{refl}$.

When normalization does not fully unify the terms (e.g.\ the loop
body contains additional logic), the path constructor applies
catamorphism fusion:

\begin{definition}[D4 Path Constructor --- formal]
\label{def:d4-path}
\[
  \mathsf{D4} \;\coloneqq\; \lambda\, f\, g\, c.\;
  \mathsf{path}\!\big(
    \mathsf{OFold}(\mathsf{op},\;
      \mathsf{OMap}(f, c),\; e),\;
    \mathsf{OFold}(\mathsf{op} \circ f,\; c,\; e)
  \big)
\]
This is the \emph{fold-map fusion} law:
$\fold(\mathsf{op}, \mathsf{map}(f, c)) = \fold(\mathsf{op} \circ f, c)$.
The auxiliary direction
$\mathsf{OMap}(f, c) = \mathsf{OFold}(\mathsf{append} \circ f,
\varepsilon, c)$ expresses the comprehension as a fold over the
input that accumulates into an initially-empty list.
\end{definition}

\begin{theorem}[D4 Soundness --- Catamorphism Fusion]
\label{thm:d4-sound}
Let $(L, \varepsilon, \mathsf{cons})$ be the initial algebra of
finite lists.  For any $f : A \to B$ and fold algebra
$(B, e, \otimes)$:
\[
  \cata_{(B,e,\otimes)} \circ \mathsf{map}(f)
  = \cata_{(B,e,\otimes \circ (\mathsf{id} \times f))}
\]
\end{theorem}

\begin{proof}
By the universal property of the initial algebra.
$\cata$ is the unique homomorphism from the initial algebra to
$(B, e, \otimes)$.  The composition $\cata \circ \mathsf{map}(f)$
is also a homomorphism from $(L, \varepsilon, \mathsf{cons})$ to
$(B, e, \otimes)$, with step $\otimes(b, f(a))$.
By uniqueness of the catamorphism, it equals
$\cata_{(B, e, \otimes \circ (\mathsf{id} \times f))}$.
\end{proof}

\textbf{Implementation mapping.}
\begin{itemize}
  \item \textbf{Normalizer}: the compiler emits \texttt{OMap} for
    both comprehensions (\texttt{\_compile\_comprehension\_ot},
    line 556) and append-loops (\texttt{\_detect\_map\_pattern},
    line 774).  The \texttt{OLam} body uses alpha-normalized
    parameter names so that $\lambda x.\,x*2$ and
    $\lambda y.\,y*2$ produce the same canonical form.
  \item \textbf{Path search}: \texttt{\_axiom\_d4\_comp\_fusion}
    (line 250) applies the fold-map fusion rule:
    when an \texttt{OFold} has an \texttt{OMap} as its collection
    and the map's filter is \texttt{None}, the map is absorbed into
    the fold by composing transforms.
\end{itemize}

\textbf{Benchmark example.}
\texttt{[x**2 for x in xs]} vs.\
\texttt{result = []; for x in xs: result.append(x**2); return result}.
Both compile to $\mathsf{OMap}(\mathsf{OLam}(x,\,
\mathsf{pow}(x, 2)),\, p_0)$.


%% ────────────────────────────────────────────────────
\section{D5: \texttt{reduce} $\leftrightarrow$ Loop}
\label{sec:d5}
%% ────────────────────────────────────────────────────

\textbf{Paradigm A} --- \texttt{functools.reduce}:
\begin{lstlisting}
def f(xs):
    return reduce(operator.add, xs, 0)
\end{lstlisting}

\textbf{Paradigm B} --- explicit accumulation loop:
\begin{lstlisting}
def g(xs):
    acc = 0
    for x in xs:
        acc += x
    return acc
\end{lstlisting}

\textbf{Compilation of A.}
\texttt{reduce(operator.add, xs, 0)} compiles to an OTerm via
Phase~3 (fold canonicalization):
$\den{f} = \mathsf{OFold}(\mathsf{iadd},\; \mathsf{0},\; p_0)$.
The normalizer recognizes \texttt{reduce} with a known operator
and rewrites it directly to a fold with the corresponding semantic
key.

\textbf{Compilation of B.}
The loop compiles to a fold via \texttt{\_detect\_augassign\_fold}:
$\den{g} = \mathsf{OFold}(\mathsf{iadd},\; \mathsf{0},\; p_0)$.
The augmented assignment \texttt{acc += x} yields semantic key
\texttt{iadd}; the init value \texttt{0} becomes the fold's initial
accumulator.

\textbf{Path.}
$\den{f} = \den{g}$ syntactically.  Path is $\mathsf{refl}$.

When the fold keys differ (e.g.\ a hash-based key from a complex
loop body vs.\ the semantic key from \texttt{reduce}), the path
constructor bridges the gap via fold universality:

\begin{definition}[D5 Path Constructor --- formal]
\label{def:d5-path}
\[
  \mathsf{D5} \;\coloneqq\; \lambda\, k_1\, k_2\, e\, c.\;
  \mathsf{path}\!\big(
    \mathsf{OFold}(k_1,\; e,\; c),\;
    \mathsf{OFold}(k_2,\; e,\; c)
  \big)
  \quad\text{when } k_1 \sim k_2
\]
where $k_1 \sim k_2$ means both keys denote the same binary
operation (verified by \texttt{\_identify\_fold\_op}).
\end{definition}

\begin{theorem}[D5 Soundness --- Fold Universality]
\label{thm:d5-sound}
Let $\otimes_1, \otimes_2 : A \times A \to A$ be binary operations.
If $\otimes_1 = \otimes_2$ as functions, and $e_1 = e_2$, then
for any collection $c$:
\[
  \mathsf{OFold}(\otimes_1, e_1, c) = \mathsf{OFold}(\otimes_2, e_2, c)
\]
\end{theorem}

\begin{proof}
By the universal property of the list catamorphism:
$\cata_{(A, e, \otimes)}$ is the unique homomorphism from the
initial list algebra to $(A, e, \otimes)$.  If $\otimes_1 = \otimes_2$
and $e_1 = e_2$, the target algebras are identical, so by uniqueness
the catamorphisms agree.
\end{proof}

\textbf{Implementation mapping.}
\begin{itemize}
  \item \textbf{Normalizer}: Phase~3 (\texttt{\_phase3\_fold},
    line 1379) rewrites builtin reducers:
    \texttt{sum(xs)} $\to$ \texttt{OFold('iadd', 0, xs)},
    \texttt{any(xs)} $\to$ \texttt{OFold('or', False, xs)},
    \texttt{all(xs)} $\to$ \texttt{OFold('and', True, xs)},
    \texttt{''.join(xs)} $\to$ \texttt{OFold('str\_concat', '', xs)}.
  \item \textbf{Path search}: \texttt{\_axiom\_d5\_fold\_universal}
    (line 276) fires on hash-named folds (8-character op\_name).
    It calls \texttt{\_identify\_fold\_op} to determine the canonical
    operation; if successful, it replaces the hash key with the
    semantic key.
  \item \textbf{HIT prover}: the OFold congruence rule (line 1274)
    handles the general case: when two \texttt{OFold}s have matching
    \texttt{init} and \texttt{collection} (checked recursively),
    it compares \texttt{op\_name}s directly, then falls back to
    \texttt{\_identify\_fold\_op} for hash-key resolution.
\end{itemize}

\textbf{Benchmark example.}
\texttt{sum(xs)} vs.\ \texttt{acc = 0; for x in xs: acc += x;
return acc}.  Phase~3 rewrites \texttt{sum} to
$\mathsf{OFold}(\mathsf{iadd}, 0, p_0)$; the loop compiles to
the same fold via \texttt{\_detect\_augassign\_fold} with key
\texttt{iadd}.


%% ────────────────────────────────────────────────────
\section{D6: Generator $\leftrightarrow$ Eager}
\label{sec:d6}
%% ────────────────────────────────────────────────────

\textbf{Paradigm A} --- generator with \texttt{yield}:
\begin{lstlisting}
def f(data):
    def helper(obj):
        if isinstance(obj, list):
            for item in obj:
                yield from helper(item)
        else:
            yield obj
    return list(helper(data))
\end{lstlisting}

\textbf{Paradigm B} --- eager list building:
\begin{lstlisting}
def g(data):
    result = []
    stack = [data]
    while stack:
        obj = stack.pop()
        if isinstance(obj, list):
            stack.extend(reversed(obj))
        else:
            result.append(obj)
    return result
\end{lstlisting}

\textbf{Compilation.}
Both programs traverse a nested structure and collect atomic elements.
The generator version yields elements lazily; the eager version
appends them to a list.  When the generator is fully consumed by
\texttt{list()}, the result is identical.

In the OTerm encoding, both compile to sequences of shared function
applications on the input.  The generator's \texttt{yield} is modeled
as producing a sequence.  Phase~4 (\texttt{\_phase4\_hof}) simplifies
$\mathsf{list}(\mathsf{gen}(\ldots))$ to the underlying sequence when
the generator is stateless.

\begin{definition}[D6 Path Constructor --- formal]
\label{def:d6-path}
\[
  \mathsf{D6} \;\coloneqq\; \lambda\, e.\;
  \mathsf{path}\!\big(
    \mathsf{OOp}(\mathsf{list},\;
      \mathsf{OOp}(\mathsf{gen},\; e)),\;
    \mathsf{OOp}(\mathsf{eager},\; e)
  \big)
\]
This is the \emph{laziness adjunction}: full consumption of a lazy
stream is naturally isomorphic to eager construction.
\end{definition}

\textbf{Path.}
$\mathsf{D6} : \Path(\mathsf{list}(\mathsf{gen}(\mathsf{body})),\;
\mathsf{eager}(\mathsf{body}))$
where $\mathsf{list} \circ \mathsf{gen}$ is the ``consume generator''
operation.  This path holds because full consumption of a generator
produces the same sequence as eager construction --- both visit the
same elements in the same order.

\textbf{Automatic inference.}
The compiler detects \texttt{yield} and \texttt{yield from} and models
the generator as producing a sequence.  The \texttt{list()} call
(or \texttt{for} loop consumption) triggers the D6 path.  The
traversal order equivalence (generator DFS vs.\ stack-with-reversed
DFS) is handled by D9 (Stack $\leftrightarrow$ Recursion).  D6 has
no dedicated axiom function because Phase~4 of the normalizer handles
the $\mathsf{list}(\mathsf{gen}(\ldots))$ pattern directly.


%% ────────────────────────────────────────────────────
\section{D7: Tail-Recursive $\leftrightarrow$ Loop}
\label{sec:d7}
%% ────────────────────────────────────────────────────

\textbf{Paradigm A} --- tail-recursive with accumulator:
\begin{lstlisting}
def f(n, acc=0):
    if n <= 0: return acc
    return f(n - 1, acc + n)
\end{lstlisting}

\textbf{Paradigm B} --- iterative with accumulator:
\begin{lstlisting}
def g(n):
    acc = 0
    while n > 0:
        acc += n
        n -= 1
    return acc
\end{lstlisting}

\textbf{Compilation.}
The tail-recursive version has the recursive call in tail position:
$f(n-1, \mathsf{acc} + n)$.  The accumulator $\mathsf{acc}$ is
passed as a parameter rather than composed in the return value
(contrast with D1 where the return value is $n \times f(n-1)$).

The compiler detects tail position and creates a RecFunction:
\[
  \mathsf{rec\_f}(n, \mathsf{acc}) = \begin{cases}
    \mathsf{acc} & n \leq 0 \\
    \mathsf{rec\_f}(n-1, \mathsf{acc} + n) & \text{otherwise}
  \end{cases}
\]

The iterative version produces the same RecFunction via the while
loop compiler (D3), with $\mathsf{acc}$ as the modified variable and
$n$ as the counter.

\begin{definition}[D7 Path Constructor --- formal]
\label{def:d7-path}
\[
  \mathsf{D7} \;\coloneqq\; \lambda\, g\, s\, \mathit{acc}_0.\;
  \mathsf{path}\!\big(
    \mathsf{OFix}(\mathsf{h},\;
      \mathsf{case}(g,\; s(\mathsf{h},\, \mathit{acc}),\;
        \mathit{acc})),\;
    \mathsf{OFold}(s,\; \mathit{acc}_0,\; \mathsf{range}(g))
  \big)
\]
This is the CPS/accumulator transform: a tail-recursive function
whose only use of the recursive result is to return it immediately
is isomorphic to a loop that threads the accumulator as mutable
state.
\end{definition}

\textbf{Path.}
$\mathsf{D7}$ identifies tail recursion with iteration.  This is the
CPS transform in action: a tail-recursive function is a loop with
the accumulator as explicit state.  Both RecFunctions have the same
defining equation, so the path is $\mathsf{refl}$.

\textbf{Automatic inference.}
The compiler detects tail position by checking whether the recursive
call is the \emph{last} operation before return (no wrapping BinOp).
When detected, it generates the same RecFunction structure as the
while-loop compiler, producing identical OTerms.
D7 is subsumed by D1 in the OTerm representation (both reduce to
the same OFix/OFold comparison), so no separate axiom function is
needed.


%% ────────────────────────────────────────────────────
\section{D8: Multi-Return $\leftrightarrow$ Single Return}
\label{sec:d8}
%% ────────────────────────────────────────────────────

\textbf{Paradigm A} --- early returns:
\begin{lstlisting}
def f(x):
    if x < 0: return -1
    if x == 0: return 0
    return 1
\end{lstlisting}

\textbf{Paradigm B} --- single return with nested conditional:
\begin{lstlisting}
def g(x):
    return -1 if x < 0 else (0 if x == 0 else 1)
\end{lstlisting}

\textbf{Compilation of A.}
The compiler produces a nested $\mathsf{OCase}$ chain:
\[
  \den{f} = \mathsf{OCase}\!\big(\mathsf{lt}(p_0, 0),\; -1,\;
    \mathsf{OCase}(\mathsf{eq}(p_0, 0),\; 0,\; 1)\big)
\]

\textbf{Compilation of B.}
The \texttt{IfExp} chain compiles to the same nested
$\mathsf{OCase}$:
\[
  \den{g} = \mathsf{OCase}\!\big(\mathsf{lt}(p_0, 0),\; -1,\;
    \mathsf{OCase}(\mathsf{eq}(p_0, 0),\; 0,\; 1)\big)
\]

When the guard ordering differs between A and B, Phase~7
(\texttt{\_phase7\_piecewise}) canonicalizes both by:
\begin{enumerate}
  \item Flattening the case chain into $(g_i, v_i)$ pieces
    via \texttt{\_flatten\_case\_chain}.
  \item Inferring the implicit else-guard from the comparison algebra
    via \texttt{\_infer\_complement\_guard}.
  \item Sorting pieces canonically by $(\mathsf{value.canon()},
    \mathsf{guard.canon()})$.
  \item Reconstructing a canonical case chain.
\end{enumerate}

\begin{definition}[D8 Path Constructor --- formal]
\label{def:d8-path}
Let $\{(g_i, v_i)\}_{i=1}^{n}$ and $\{(g'_{\sigma(i)},
v'_{\sigma(i)})\}_{i=1}^{n}$ be two enumerations of the same
piecewise function (i.e.\ the guards form a partition and
$v_i = v'_{\sigma(i)}$ on the region $g_i$).  Then:
\[
  \mathsf{D8} \;\coloneqq\; \lambda\, \vec{g}\, \vec{v}\, \sigma.\;
  \mathsf{path}\!\big(
    \mathsf{OCase}(g_1, v_1, \mathsf{OCase}(g_2, v_2, \ldots)),\;
    \mathsf{OCase}(g'_{\sigma(1)}, v'_{\sigma(1)},
      \mathsf{OCase}(g'_{\sigma(2)}, v'_{\sigma(2)}, \ldots))
  \big)
\]
This is the \emph{section merging} path: two case chains over the
same partition are equal as piecewise functions.
\end{definition}

\begin{theorem}[D8 Soundness --- Section Merging]
\label{thm:d8-sound}
Let $\{U_i\}_{i=1}^{n}$ be a finite cover of the input space by
decidable predicates (guards).  Let $s_i, s'_i : U_i \to A$ be
sections such that $s_i = s'_{\sigma(i)}$ on each $U_i$.  Then
the piecewise functions $\bigsqcup_i s_i$ and
$\bigsqcup_i s'_{\sigma(i)}$ are equal.
\end{theorem}

\begin{proof}
This is the sheaf condition for the constant presheaf on a finite
decidable cover.  The sections agree on every fiber of the cover, and
since the cover is a partition (the guards are disjoint and
exhaustive), there is a unique global section.  Reordering the
enumeration of a partition does not change the global section.
Formally, this follows from the fact that $\Hone = 0$ for the
constant presheaf on a finite cover by decidable predicates.
\end{proof}

\textbf{Implementation mapping.}
\begin{itemize}
  \item \textbf{Normalizer}: Phase~7 (\texttt{\_phase7\_piecewise},
    line 1696) is the primary mechanism.  It calls
    \texttt{\_try\_canonicalize\_piecewise} (line 1742) which
    flattens, infers complement guards, sorts, and reconstructs.
  \item \textbf{Path search}: \texttt{\_axiom\_d8\_section\_merge}
    (line 313) handles residual nested-case flattening after other
    axioms have fired.  It detects
    $\mathsf{OCase}(g, A, \mathsf{OCase}(g', B, C))$ where $g$
    and $g'$ are complementary guards (checked by
    \texttt{\_guards\_complementary}), and flattens to
    $\mathsf{OCase}(g, A, C)$.
  \item \textbf{HIT prover}: \texttt{\_try\_case\_flatten}
    (line 1552) handles the cross-case permutation:
    $\mathsf{case}(t_1, A, \mathsf{case}(t_2, B, C))$ vs.\
    $\mathsf{case}(t_2, B', \mathsf{case}(t_1, A', C'))$.
    It verifies that the guards and branches match up to
    reordering, proving equivalence of two enumerations of the
    same piecewise function.
\end{itemize}

\textbf{Benchmark example.}
The sign function: one version uses two \texttt{if} guards with
early returns (negative, zero) and a fallthrough positive case;
the other uses a single ternary expression.  Phase~7 flattens both
to the canonical piecewise form
$\mathsf{OCase}(\mathsf{eq}(p_0, 0), 0,
\mathsf{OCase}(\mathsf{lt}(p_0, 0), -1, 1))$
(sorted by value then guard), producing $\mathsf{refl}$.


